\subsection{Global properties of gene-trait relationships}

"Our examples show that the global map can recapitulate manual searching. Its real power though comes from it being complete and quantitative. This motivates these landscape analyses..."

Plot of sum of betas across traits
Plot of number of genes with high total vs. number with high indirect (points are traits); also number of genes with high indirect vs number with high direct
How many greater than 50\%? How many unique genes of these?
How many greater than 50\% on average for a trait? Use 5\% threshold for common and 0.5\% threshold for rare

Make plot of sum direct vs sum indirect.
 - the sum of indirect is in general larger
 - common traits have correlation more strongly with pLI than rare (is this a power issue though)
 - pLI correlates with sum of indirect and sum of direct, but direct more strongly (confounders in trait ascertainment, confounders in power may be an issue here)
  - for genes like APOC4, a lot of traits with no indirect. But how much of that is due to lack of power of the GWAS? We can't say necessarily that those traits are novel findings for it


Amount of redundant information in gene set libraries -- plot sum of betas vs. sum of betas uncorrected. The difference indicates how many gene sets are actually redundant.

Sum of betas for rare vs. sum of betas for common -- any more or less coherence 

Plot number of clumps vs sum of betas

Plot sum indirect or sum direct as a function of number of annotations (have to filter and correct for traits first)

How many effector gene predictions do we make
 - how many cases does probability without gene sets increase after with gene sets
 - how many genes near a lead SNP have probability above X

 Histogram of number of effective genes per loci (need the ge.out)

Investigate what lack of direct means for rare. pLoF/LOEUF vs. indirect-direct etc. The ones with indirect + direct likely more constrained, but indirect in general more constrained than non-indirect

Then what lack of indirect means for GWAS. Maybe here we have histogram of number of genes per locus with prior above X. Show what happens to direct with gene sets vs. direct w/o gene sets

\subsection{Calibrating non-human-genetic data sources across different human traits}

\begin{figure}[t!]
  \centering
  \includegraphics[width=1\linewidth]{figures/indirect-support/real/validation-figure-1}
  \caption{(a) A subset of pigean models run on a 13 trait common disease validation set. Displaying top performing models and models selected through assessment of top genesets. Marker size represents the complexity of the model and the multiple transparent grey boxes show overlayed null distrbutions in terms of both relative success and mechanistic alignment for each model (Methods). (b) Same models run over a 40 trait rare disease valiation set. (c) Over all model selection experiments the proportion of genesets that came from each geneset library that where significant for a given trait-model pair. }
  \label{fig:model-selection-main}
\end{figure}

Which libraries are good/bad
Which libraries vary in utility across trait groups